\documentclass[10pt,a4paper]{amsart}
\usepackage[margin=19mm]{geometry}
\usepackage{amsmath,amssymb,mathtools}
\usepackage[scr=rsfs,cal=euler]{mathalfa}
\usepackage{xfrac}
\usepackage[unicode,hidelinks,bookmarks=false]{hyperref}
\newcommand{\ie}{i.e.\ }
\newcommand{\cf}{cf.\ }
\newcommand{\resp}{respectively\ }
\newcommand{\Subsection}{\subsection}
\providecommand{\nobreakdash}{\nobreak\hbox{-}}
\setlength{\emergencystretch}{2em}
\multlinegap0pt
\usepackage{bm}
\usepackage{bbm}
\usepackage{dsfont}
\usepackage{xspace}
\usepackage{cancel}
\usepackage{mathtools}
\usepackage{amsthm}
\makeatletter
\@ifundefined{teo}{\newtheorem{teo}{Teo}}{}
\@ifundefined{defin}{\newtheorem{defin}[teo]{Definition}}{}
\@ifundefined{prop}{\newtheorem{prop}[teo]{Proposition}}{}
\makeatother
\def\R{\mathbb{R}}
\title[Critical singular problems in Carnot groups]{Critical singular problems in Carnot groups}
\author{Stefano Biagi, Mattia Galeotti, Eugenio Vecchi}
\date{Source: arXiv:2506.07521v1 (CC BY 4.0)}
\begin{document}
\maketitle

\noindent\emph{Source and licence.} Critical singular problems in Carnot groups. Version arXiv:2506.07521v1 (CC BY 4.0), distributed under the Creative Commons Attribution 4.0 International licence, as recorded in the arXiv licence metadata for this exact version and shown on the abstract page https://arxiv.org/abs/2506.07521v1. Full rights notice: licenses/arxiv-2506.07521-CC-BY-4.0.txt.\par\smallskip

\noindent\textit{Editorial scope.} Counted unit: the Prop whose source label is \texttt{prop:SMP} in the article (source0.tex lines 674--682, source label \texttt{prop:SMP}), delivered together with the article's own complete original proof of it (source0.tex lines 683--745, one \texttt{proof} environment of 63 lines). No mathematics is written, repaired or re-derived: statement and proof are the article's own text line for line. Required context retained uncounted: prop:WHarnack (lines 609--629); eq:PDEOrderzeroHarnack (lines 614--616); eq:General\_Singular\_Problem (lines 749--755); eq:Growth\_of\_f (lines 757--759); def:weak\_sub\_super\_sol (lines 762--772). Every definition, display and label of those lines is retained unchanged. Omitted from this package and not redistributed: 1--608, 617--673, 746--748, 756--756, 760--761, 773--3683. Nothing inside a retained excerpt is omitted or altered: every retained body is delivered exactly as the article's lines are written, so no omission receipt of any kind accompanies this package. Citations: the retained excerpts carry no citation command at all, so no bibliography entry of the article is transcribed and no reference is redistributed. Reference adaptation: none was needed and none was made; every reference inside the delivered unit resolves inside it, so no unresolved reference or citation is produced. Printed slips kept literal: no printed word, symbol, hypothesis or number of the counted statement or of its proof was altered. Layout and class: the article's own class is \texttt{amsart} with packages including amsfonts, enumerate, amssymb, graphicx, graphics, stmaryrd, geometry, babel, inputenc, amsthm, amsmath, amstext, tipa, multirow; the wrapper is the lane's pinned \texttt{amsart} editorial wrapper at 10pt on A4, which restates the theorem-like environments the retained text needs and adds the article's own macro definitions that the retained text needs (\texttt{\textbackslash R}), with extra wrapper packages (bm, bbm, dsfont, xspace, cancel, mathtools). The delivered numbering is the unit's own. Assets: no retained span contains a figure environment, a graphics-inclusion command, a PDF-page inclusion, a table or a TikZ picture; no image file is copied into this package, the package declares no asset dependency and no image file is redistributed with it. Licence: version arXiv:2506.07521v1 of this article is distributed under the Creative Commons Attribution 4.0 International licence, as recorded in the arXiv licence metadata for this exact version and shown on the abstract page https://arxiv.org/abs/2506.07521v1. A full rights notice is delivered at licenses/arxiv-2506.07521-CC-BY-4.0.txt. Characters: every retained excerpt is pure ASCII, so the delivered engine, XeLaTeX with its default Latin Modern text font, reports no missing character.
\begin{prop}[Weak Harnack inequality for $-\Delta_\mathbb{G}+c$] \label{prop:WHarnack}
 Let $\mathcal{O}\subseteq\R^n$ be a bounded open set, and let 
 $c\in L^\infty(\mathcal{O}),\,c\geq 0.$
 Moreover, let $u\in S_0^1(\mathcal{O})$ be a \emph{weak supersolution}
 of the equation 
 \begin{equation} \label{eq:PDEOrderzeroHarnack}
  \text{$-\Delta_\mathbb{G}u+cu = 0$ in $\mathcal{O}$},
  \end{equation}
 that is
 \emph{(}see also the subsequent Definition \ref{def:weak_sub_super_sol}\emph{)},
 $$\int_\mathcal{O}\langle\nabla_\mathbb{G}u,\nabla_\mathbb{G}\varphi\rangle_{\mathfrak{g}_{1}}
 +\int_\mathcal{O} cu\varphi\geq 0\quad\text{
 for every $\varphi\in C_0^\infty(\mathcal{O}),\,\varphi\geq 0$}.$$
 If, in addiction, $u\geq 0$ a.e.\,in $\mathcal{O}$, there exist
 constants $c_0 > 0$ and $p_0\in (0,1)$,
 \emph{independent of~$u$}, such that,
 for every $g_0\in\mathcal{O},\,r > 0$ with $B_{4r}(g_0)\Subset\mathcal{O}$, we have
 \begin{equation} \label{eq:HarnackInequ}
  \Big(-\!\!\!\!\!\!\!\int_{B_{3r}(g_0)}|u|^{p_0}\Big)^{1/p_0}\leq c_0\,\inf_{B_{r}}u.
 \end{equation}
\end{prop}

\begin{equation}\label{eq:General_Singular_Problem}
	\left\{\begin{array}{rl}
		-\Delta_{\mathbb{G}}u = \dfrac{\lambda}{u^{\gamma}} + f(x,u)& \textrm{ in } \Omega,\\
		u>0 & \textrm{ in } \Omega,\\
		u=0 & \textrm{ on } \partial \Omega,
	\end{array}\right.
\end{equation}

\begin{equation}\label{eq:Growth_of_f}
	|f(x,t)|\leq K_{f} (1+t^{2^{\star}_{Q}-1})  \quad \textrm{ for a.e. } x \in \Omega \textrm{ and for  every } t>0.
\end{equation}

\begin{defin}\label{def:weak_sub_super_sol}
	Let $f: \Omega \times (0,+\infty)\to \mathbb{R}$ be a Carath\'{e}odory function satisfying the growth condition \eqref{eq:Growth_of_f}. We say that a function $u \in S^{1}_{0}(\Omega)$ is a weak subsolution (resp. supersolution) of \eqref{eq:General_Singular_Problem} if it satisfies the following properties:
	\begin{itemize}
		\item[(i)] $u>0$ in $\Omega$ and $u^{-\gamma} \in L^{1}_{\textrm{loc}}(\Omega)$.
		\item[(ii)] For every $0\leq \varphi \in C^{\infty}_{0}(\Omega)$, it holds that
		\begin{equation}\label{eq:weak_sub_super_sol}
			\int_{\Omega}\langle \nabla_{\mathbb{G}}u, \nabla_{\mathbb{G}}\varphi \rangle_{\mathfrak{g}_{1}} \leq (\textrm{resp. } \geq) \lambda \int_{\Omega}u^{-\gamma}\varphi + \int_{\Omega}f(x,u)\varphi.
		\end{equation}
	\end{itemize}
	Finally, we say that $u \in S^{1}_{0}(\Omega)$ is a weak solution of \eqref{eq:General_Singular_Problem} if it is both a weak subsolution and a weak supersolution of \eqref{eq:General_Singular_Problem} without the non-negativity condition on $\varphi$.
\end{defin}

\begin{prop}[Strong Maximum Principle for $L$] \label{prop:SMP}
 Let $\mathcal{O}\subseteq \mathbb{G}$ be a \emph{bounded and connected} 
 open set, and let
 $u\in S^1_0(\mathcal{O}),\,u\geq 0$ be a \emph{weak supersolution} of \eqref{eq:PDEOrderzeroHarnack}. 
 Then,
 \begin{equation} \label{eq:SMPConclusion}
  \text{either $u \equiv 0$ or $u > 0$ a.e.\,in $\mathcal{O}$}.
 \end{equation}
\end{prop}

\begin{proof} 
 We assume that there exists a set $\mathcal{Z}\subseteq\mathcal{O}$ of \emph{positive Lebesgue measure} such that $u \equiv 0$ (pointwise) on $\mathcal{Z}$, and we prove that in this case we have $u\equiv 0$ a.e.\,on $
\mathcal{O}$.
\vspace{0.1cm}

To this end we first observe that, since $|\mathcal{Z}| > 0$, we can find $r_0 > 0$ such that
$$\mathcal{Z}_0 = \mathcal{Z}\cap \{g\in\mathcal{O}:\,d_\mathbb{G}(g,\partial\mathcal{O}) \geq r_0\}$$
 has \emph{positive Lebesgue measure}; moreover, since $K = \{g\in\mathcal{O}:\,d_\mathbb{G}(g,\partial\mathcal{O}) \geq r_0\}\subseteq\Omega$
 is  com\-pact (recall that $d_\mathbb{G}$ induces 
 the Euclidean topology), there exist $g_1,\ldots,g_p\in K$ and $r_1,\ldots,r_p > 0$ (for some $p\in\mathbb{N}$) such that $B_{4r_i}(g_i)\Subset\mathcal{O}$ (for $1\leq i\leq p$) and
 $$\mathcal{Z}_0 = \bigcup_{i = 1}^p (\mathcal{Z}_0\cap B_{r_i}(g_i)).$$
 As a consequence, there exists a $d_\mathbb{G}$-ball $B_0 = B_{r_{i_0}}(g_{i_0})$
 (for some $1\leq i_0\leq p$) such that
 \begin{equation} \label{eq:Z0capBpositiveMeas}
  |\mathcal{Z}_0\cap B_{r_{i_0}(g_{i_0})}| > 0.
 \end{equation}	
 In particular, since $u\geq 0$ a.e.\,in $\mathcal{O}$ and $u = 0$ on $\mathcal{Z}_0$, we have
 \begin{equation} \label{eq:infuzeroZ0}
 	\inf_{B_{r_{i_0}}(g_{i_0})} u = \inf_{\Omega}u = 0.
 \end{equation}
 With \eqref{eq:Z0capBpositiveMeas}-\eqref{eq:infuzeroZ0}
at hand, we then define the set
 $$S = \{g\in\mathcal{O}:\,\text{there exists a ball 
 $B = B_r(g)\subseteq\mathcal{O}$ 
 s.t.\,$u = 0$ a.e.\,on $B$}\},$$
 and we prove that $S = \mathcal{O}$ by a \emph{connection argument}:
  \vspace{0.1cm}
 
 -\,\,\emph{$S\neq\emptyset$}. Let $g_{i_0}\in\mathcal{O}$ be as in \eqref{eq:Z0capBpositiveMeas}-\eqref{eq:infuzeroZ0}. Since $u$ is a non-negative weak supersolution of equation \eqref{eq:PDEOrderzeroHarnack}, and since
 $B = B_{4r_{i_0}(g_{i_0})}\Subset\mathcal{O}$, we can apply the Weak Harnack Inequality in Proposition \ref{prop:WHarnack}: this gives the following estimate
 $$\Big(-\!\!\!\!\!\!\!\int_{B_{3r_{i_0}}(g_{i_0})}|u|^{p_0}\Big)^{1/p_0}\leq c\,\inf_{B_{r_{i_0}}(g_{i_0})}u = 0,
$$
from which we derive that $u = 0$ a.e.\,in $B_{3r_{i_0}}$. Hence, $g_{i_0}\in S$.
 \medskip
 
 
 -\,\,\emph{$S$ is open}. Assume that $g\in S$, and let $B = B_r(g)\subseteq \mathcal{O}$ be a 
 $d_\mathbb{G}$-ball such that  
 $$\text{$u = 0$ a.e.\,in $B$}.$$ 
 Given any $g'\in B$, if we choose $\rho > 0$ in such a way that $B_\rho(g')\subseteq B$, we obviously have
  $$\text{$u = 0$ a.e.\,on $B_{\rho}(g')\subseteq B$},$$
 and this proves that $g'\in S$.
 By the arbitrariness of $g'\in B_r(g)$, we  conclude that $B_r(g)\subseteq S$, and hence $S$ is open, as desired.

 \medskip

-\,\,\emph{$S$ is closed}. Assume that $(g_k)_k$ is a sequence of points in $S$ which converges (as $k\to+\infty$) to some $g\in\mathcal{O}$, and let $\rho > 0$
be such that $B_{4\rho}(g)\Subset\mathcal{O}$. Since $g_k\to g$ as $k\to+\infty$, there
exists some $k_0\in\mathbb{N}$ such that $g_{k_0}\in B_\rho(g)$; on the other hand, since
$g_{k_0}\in S$, it is possible to find a suitable $d_\mathbb{G}$-ball $B = B_r(g_{k_0})$
such that $B \subseteq B_\rho(g)$ and
$$\text{$u = 0$ a.e.\,in $B$}.$$
Hence, using once again the 
Weak
 Harnack Inequality in Proposition \ref{prop:WHarnack}, we get
$$\Big(-\!\!\!\!\!\!\!\int_{B_{3\rho}(g)}|u|^{p_0}\Big)^{1/p_0}\leq c\,\inf_{B_{\rho}(g)}u = 0,
$$
and thus $u = 0$ a.e.\,in $B_{3\rho}(g)$, that is, $g\in S$. Hence, $S$ is closed.
\medskip

Gathering the above facts, and recalling that $\mathcal{O}$ is connected, we then conclude that $\mathcal{O} = S$, and this obviously implies that
$u = 0$ a.e.\,in $\mathcal{O}$, as desired.
\end{proof}

\end{document}
